\section{Introduction}
\label{sec:introduction}

Clusterized quantum magnets provide a direct route from the spectrum of a
few-spin unit to the collective behavior of a quantum many-body system.
Their local multiplets can survive as low-energy degrees of freedom, become
dispersive, or be removed from an experimental response by symmetry.
Clusterized chains therefore separate three questions that are often mixed
together in a conventional phase diagram: which local state remains active, how
neighboring units couple, and which resulting excitations are visible.

The spin-$1/2$ tetramer chain is a well-studied reference of this type, because
spin order, entanglement structure, and excitation spectra have been organized
within a common coupling diagram~\cite{li2026tetramer,gong2008tetramer}.
Experimentally, the tetramer compound CuInVO$_5$ has been characterized as a
multimerized magnet with a singlet unit-cell ground state, and its magnetic
spectrum has been discussed in terms of both local multiplet and spinon
excitations~\cite{hase2016cuinvo5,reja2019cuinvo5,singhania2018}.
Trimer chains, by contrast, host a doublet unit cell and support fractionalized
and composite excitations that have been resolved in
RIXS~\cite{cheng2022trimer,cheng2024trimerfield,li2025rixstrimer,prabhakar2025trimer,bera2022trimer}.
Beyond the spectrum, the dimerized and tetramerized limits are classified by
nonlocal string order and symmetry-protected
topology~\cite{guwen2009,pollmann2010,pollmann2012,han2017bcfg,ma2023aspt},
diagnosed by order parameters and entanglement measures that do not rely on a
local multiplet picture~\cite{hung2015,lu2020,scheie2021}.
Related one-dimensional models further show that an intermediate deconfined
critical state can separate two gapped cluster
phases~\cite{jiang2019dqcp,roberts2019dqcp,huang2019dqcp,yang2021dqcp}, a
possibility that has to be tested explicitly for an odd unit cell.
Even in the molecular limit, the collective spin dynamics of a closed cluster is
not simply a sum of local transitions~\cite{ummethum2012}.

Replacing four spins by five is not a minor geometric modification.
A unit cell with an odd number of half-integer spins retains a half-integer
degree of freedom unless symmetry breaking or ground-state degeneracy
intervenes.
Lieb--Schultz--Mattis-type results therefore rule out a unique
symmetry-preserving gapped ground state under the relevant translation and
spin-rotation assumptions~\cite{lieb1961,ogata2019,affleck1987}, in close analogy
with the Haldane conjecture for integer-spin chains~\cite{haldane1983} and the
dimerized-chain crossovers studied earlier~\cite{hida1992,chitra1995}.
The odd cell thus changes the question.
It is no longer which gapped cluster phase forms, but which cell state carries the
low-energy physics, whether the resulting chain is the natural $c=1$ Luttinger
liquid~\cite{giamarchi2003,calabrese2004}, and which of its excitations a spin
probe can reach.

Here we study a linear Heisenberg pentamer chain with two intra-cell exchanges,
$J_1$ and $J_2$, and an inter-cell exchange $J_3$.
We set $J_1=1$ and define $\alpha=J_2/J_1$ and $\beta=J_3/J_1$.
The geometry differs from previously studied mixed-spin or differently connected
pentamers~\cite{haraldsen2011pentamer}.
The isolated cell is the natural starting point: its ground multiplet fixes the
effective degree of freedom of the chain, and its reflection parity fixes which
local excitations a spin probe can excite.
We therefore determine the complete isolated-cell spectrum, construct the
projected low-energy chain, and test this picture against static and dynamical
many-body calculations.

Three results organize this work.
First, throughout the sampled positive-$\alpha$ region, the pentamer ground
state is an even-parity spin-$1/2$ doublet and positive $\beta$ generates an
antiferromagnetic effective chain.
Second, reflection parity makes the lowest excited doublet essentially dark at
$q=\pi$ and concentrates the local inelastic weight in spin-$3/2$ multiplets, in
parallel with the momentum-selective response discussed for related clusterized
chains~\cite{gao2024srv3o9,schlappa2018,ishii2025kedge}.
Third, negative $\alpha$ switches the isolated cell to a spin-$3/2$ quartet whose
folded shape factor peaks at $0.60\pi$, and finite chains develop a momentum
reconstruction from that value toward $\pi$.
Two independent diagnostics of this reconstruction select different coupling
intervals, so it is best described as a finite-size crossover rather than a phase
boundary.
These results rest on the standard methodology of one-dimensional magnetism: matrix-product-state
DMRG~\cite{white1992,schollwoeck2005,schollwoeck2011} combined with
correction-vector and real-time spectral
techniques~\cite{kuehner1999,jeckelmann2002,vidal2004,paeckel2019}, supported by
finite-temperature and entanglement-based
diagnostics~\cite{barthel2009,fishman2022,keselman2020}.

Section~\ref{sec:methods} defines the model, the projection, and the numerical
criteria.
Section~\ref{sec:results} follows the positive-coupling chain from the isolated
cell to the collective spectrum and then turns to negative $\alpha$.
Section~\ref{sec:discussion} compares the results with the tetramer chain and
states their consequences for scattering experiments.
